\documentclass{swfulabjournal}

\swfusetup{%
  % Title ={《计算机网络》课程实习}, % 课程名称
  Title   ={综合实习二}, % 课程名称
  Author  ={}, % 
  ID      ={20231159000}, % 学号
  Subject ={计算机科学与技术},%
  Grade   ={2023级},%
  Tutor   ={王晓林}, %指导教师
  YearS   ={\the\year},%
  YearE   ={\the\year},%
  MonthS  ={6},%
  MonthE  ={6},%
  DateS   ={28},%
  DateE   ={30},%
}

\begin{document}

\maketitle % 封面

\addjournal%
{%实习日期
  2024年6月28日}%
{%实习时间
  8:20 --- 18:20}%
{%实习地点
  经管楼219}%
{%实习单位
  校本部}%
{%实习内容
  熟悉实验环境，学习Linux命令行的使用，查阅相关网络命令的手册。}%
{%收获与体会
  掌握了常用Linux命令的使用，了解了常用网络命令的功用和操作方法。}%
\clearpage
\addjournal%
{%实习日期
  2024年6月29日}%
{%实习时间
  8:20 --- 18:20}%
{%实习地点
  经管楼219}%
{%实习单位
  校本部}%
{%实习内容
  学习在Linux平台抓取网络数据包，并进行逐字节分析}%
{%收获与体会
  了解了如何在命令行利用nc、tcpdump等网络工具搭建实验环境，并抓取、分
  析数据包。}%
\clearpage
\addjournal%
{%实习日期
  2024年6月30日}%
{%实习时间
  8:20 --- 18:20}%
{%实习地点
  经管楼219}%
{%实习单位
  校本部}%
{%实习内容
  学习使用ttyrec和tmux进行命令行操作的录屏，并撰写实习报告。}%
{%收获与体会
  学会了利用ttyrec和tmux录制HTTP实验的过程。}%

\end{document}
%%% Local Variables:
%%% mode: LaTeX
%%% TeX-master: t
%%% End:
